\section{A conditional improvement using exact prime removal}
\label{sec:qrh-bootstrap}

The absolute deletion error in Lemma~\ref{lem:qrh-amplification}
loses information about the M\"obius function. Retaining that
information gives a better conditional implication. The new
mean-square estimate required by the implication is stated explicitly;
it is not established by the supplied estimates in their present ranges.

Use the notation of Lemma~\ref{lem:qrh-amplification}, fix
$k\ge2$, and suppose that the finite-order twists and their zero
extensions satisfy
\[
 A_1(T)=A_{p^k}(T)-\nu(p)A_{p^k}(T/Np)
       \qquad(p\notin S).
\]
For a fixed compactly supported $W$, all these sums vanish for
sufficiently small positive $T$. This point makes the following
iteration exact, without an infinite unestimated remainder.

\begin{lemma}[Uniform control of the deleted prime]
\label{lem:qrh-better-deletion}
Suppose $\beta>0$ and
$A_1(T)\ll_\epsilon T^{\beta+\epsilon}$ for every
$\epsilon>0$, with fixed arithmetic data and fixed $W$.
Then, uniformly in primes $p\notin S$ and all $T>0$,
\[
 A_{p^k}(T)=\sum_{j\ge0}\nu(p)^j A_1(T/(Np)^j),
 \qquad
 |A_{p^k}(T)|\ll_\epsilon T^{\beta+\epsilon}.
\]
Consequently, for $Y/2<Np\le Y$,
\[
 |A_1(D)-A_{p^k}(D)|
       \ll_\epsilon(D/Y)^{\beta+\epsilon}.
\]
\end{lemma}
\begin{proof}
Solve the displayed prime-removal identity for $A_{p^k}(T)$
and iterate. For each $T>0$, the sum $A_{p^k}(T/(Np)^J)$
vanishes once $J$ is large. Thus the expansion is a finite identity.
Choose $v>0$ containing the upper endpoint of the support
of $W$. The assumed estimate can be made valid for every
$T>0$ by increasing its constant: the sum is zero for
$T<1/v$, and is bounded on every fixed compact range of
positive $T$. Hence
\[
 |A_{p^k}(T)|\le
 C_\epsilon T^{\beta+\epsilon}
       \sum_{j\ge0}(Np)^{-j(\beta+\epsilon)}
 \le\frac{C_\epsilon T^{\beta+\epsilon}}
                {1-2^{-\beta-\epsilon}}.
\]
The denominator is positive and independent of $p$.
Apply this bound at $T=D/Np$, and use $Np>Y/2$,
to obtain the last assertion.
\end{proof}

\begin{proposition}[A conditional exponent improvement]
\label{prop:qrh-exponent-improvement}
Let $1/2<\beta\le1$ and $0<h<k$. Suppose that, for every
$\epsilon>0$, the following two estimates hold for the same
fixed $W$ and arithmetic data:
\[
 A_1(D)\ll_\epsilon D^{\beta+\epsilon},\qquad
 \sum_{0<Nu\le D^h}|A_u(D)|^2
       \ll_{h,\epsilon}D^{1+h+\epsilon}.
\]
Suppose also that the prime-ideal count used in
Lemma~\ref{lem:qrh-amplification} is available. Then
\[
 A_1(D)\ll_{h,\epsilon}D^{\beta'(h)+\epsilon},
 \qquad
 \beta'(h)=\max\left\{
    \frac12+\frac{k-1}{2k}h,
    \beta\left(1-\frac hk\right)\right\}.
\]
In particular, if the mean-square estimate is available at
\[
 h=h_*(\beta):=\frac{k(2\beta-1)}{k-1+2\beta},
\]
then
\begin{equation}\label{eq:qrh-bootstrap-map}
 \beta'=T_k(\beta):=\frac{k\beta}{k-1+2\beta}<\beta.
\end{equation}
All implied constants may depend on $\beta,h,W$ and on the
fixed character, but must be independent of the varying prime
$p$ and of $D$.
\end{proposition}
\begin{proof}
Set $H=D^h$ and $Y=H^{1/k}$. The proof of
Lemma~\ref{lem:qrh-amplification}, now using
Lemma~\ref{lem:qrh-better-deletion}, gives
\[
 |A_1(D)|^2\ll
 D^{1+h-h/k+\epsilon_1}\log(2D)
       +(D/Y)^{2(\beta+\epsilon_2)}.
\]
Choose the preliminary positive losses $\epsilon_1,\epsilon_2$
small enough in terms of any requested final loss, absorb
$\log(2D)$ into a further arbitrarily small power, and take
square roots. The first resulting exponent increases with
$h$ and the second decreases. They coincide at $h_*(\beta)$;
substitution gives \eqref{eq:qrh-bootstrap-map}. Since
$\beta>1/2$, one has $0<h_*(\beta)<k$ and
$k-1+2\beta>k$, so the new exponent is strictly smaller.
\end{proof}

For sixth-power characters this would give
\[
 \beta=\frac78,\quad h_* =\frac23,
       \quad T_6(\beta)=\frac79;
 \qquad
 \beta=\frac{11}{12},\quad h_* =\frac{30}{41},
       \quad T_6(\beta)=\frac{33}{41}.
\]
Thus the first useful extension of the known $7/8$ result
would concern a mean square over $Nu\le D^{2/3}$.
That is substantially shorter than the available range
$Nu\le D^{1+\eta}$.

If the hypothesis at $h_*(\beta)$ could be deduced from the
current exponent bound for every $\beta>1/2$, iteration
would approach the critical line. More precisely, if
$\beta_{n+1}=T_k(\beta_n)$, then
\[
 \frac1{\beta_n}
 =2-\left(\frac{k-1}{k}\right)^n
               \left(2-\frac1{\beta_0}\right),
 \qquad \beta_n\downarrow\frac12.
\]
This follows by taking reciprocals in
\eqref{eq:qrh-bootstrap-map} and solving the affine
recurrence. For any fixed zero to the right of $1/2$,
only finitely many iterations would be needed to obtain
a contradictory bound in Proposition~\ref{prop:qrh-mellin}.
Thus constants may deteriorate with the number of
iterations; no uniform limit of those constants is being
assumed. For the conclusion about every Hecke character, these
implications must hold for every fixed character and every
compactly supported smooth $W$, with a common exponent;
their implied constants may depend on those fixed data.
The missing assertion is the implication from the
current exponent to the shorter mean square.

\subsection{What a fixed lower limit on the range permits}

Suppose the mean-square estimate is available only at
exponents $h\ge h_0$, where $0<h_0<k$. The best consequence of the last
proposition is obtained at $h=\max(h_0,h_*(\beta))$
whenever this lies in the permitted range. In particular,
a strict improvement over $\beta$ requires
\[
 h_0<\frac{k(2\beta-1)}{k-1}.
\]
Indeed the first entry in $\beta'(h)$ must be smaller than
$\beta$, while the second is smaller for every $h>0$.
A fixed $h_0$ leaves the lower bound
$1/2+(k-1)h_0/(2k)$ on the exponent obtainable by this
particular argument. With $k=6$ and $h_0=1$, that bound
is $11/12$. A previously known $7/8$ estimate therefore
does not improve when inserted into the existing
$h>1$ mean-square range.

One must also avoid a circular reformulation. If the bound
\[
 \sum_{0<Nu\le D^h}|A_u(D)|^2\ll_{h,\epsilon}
                         D^{1+h+\epsilon}
\]
were simply assumed for every $h>0$, its single term
$u=1$ would already give
$A_1(D)\ll_\epsilon D^{1/2+\epsilon}$ by taking $h$
sufficiently small. Amplification would not be needed.
Even at $h=h_*(\beta)$, the single term gives a strict,
though weaker, improvement, because
$h_*(\beta)<2\beta-1$. Thus proving this shorter
mean-square estimate is itself substantive new cancellation;
it is not a formal consequence of changing parameters.

\subsection{The precise failure of the existing iteration in this range}

The October proof transforms the original scale $D$ and
range $H$ by Poisson summation. In the contribution whose
common-divisor parameter is $1$, its new parameters have
sizes
\[
 X=D,\qquad F=1,\qquad \mathcal H\asymp D^2/H,
          \qquad\Sigma=XF=D.
\]
The mean-square theorem in its Proposition 5.1 requires
$\mathcal H\le\Sigma D^{-\kappa}$ for a fixed
$\kappa>0$. This requires $H$ to be greater than $D$
by a fixed power. If $H=D^h$ with $h<1$, the ratio
$\mathcal H/\Sigma\asymp D^{1-h}$ instead tends to
infinity. Fixed scale constants do not change this
conclusion.

This is also visible directly in the cube-divisor step.
The splitting threshold in its equation (5.6) satisfies
\[
 H_c^3=\min(X,X^2/\mathcal H^2)<1
\]
for large $D$. The unit divisor $b=1$ is therefore in
the remaining sum, and its new length is $L_b=X$.
An attempted reuse of the displayed supremum for that
remaining sum would include the original mean square itself.
There is no decrease of the parameters. In fact, the lemma
stating that bound already assumes $\mathcal H\le\Sigma$;
we are diagnosing why its proof does not extend to this
case, not applying the lemma outside its hypotheses. The next transfer proposition requires
$\mathcal H\le\Sigma$ and is unavailable at the desired
$\Sigma=X$. Thus the existing finite induction does not
prove the shorter-range estimate.

In the September paper, the unmarked case of the principal
inverse moment requires $r<m-c$ when the norm length is
$D=Z^r$ and the outer range is $H=Z^m$. Its canonical
estimate also requires, even when the additional excluded
ideal is the unit ideal,
$1\le Z^{F_0-M-c}$ with $c>0$. These conditions enforce
the same strict direction of inequality between the
corresponding ranges. The longer-length estimate later
obtained for sixth-power-free elements does not remedy this
problem: the amplifying elements $p^6$ have valuation six
and are excluded from that index set.

These are exact limitations of the statements and
recurrences presently available. They do not prove that a
new argument cannot extend the mean square. They locate
where such an argument is needed and prevent a purported
iteration to the critical line from concealing that step.

\subsection{A counterexample to enlarging the coefficient class}

The shorter mean square also cannot be obtained by applying an
identical estimate for arbitrary bounded squarefree coefficients.
Here is an explicit obstruction. Fix $0<h<k/(k-1)$, put
$H=D^h$, and take the fixed column coefficient $c_{\mathfrak a}$
to be $1$ when $\mathfrak a$ is prime outside $S$, and
zero otherwise. This coefficient sequence is independent
of $D$ and of the outer index. Choose a fixed nonnegative
$W\in C_c^\infty((1/3,3/2))$ which is at least one
on $[1/2,1]$. For every prime ideal $p\notin S$ of norm in
$[H^{1/k}/2,H^{1/k}]$, one has $Np<D/3$ when $D$ is
large, because $h/k<1$. Thus every prime column on which
$W(N\mathfrak a/D)\ne0$ is different from $p$. Thus the zero-extended character
is exactly $1$ on every column at the argument $p^k$.
The prime-ideal count gives
\[
 \left|\sum_{\mathfrak a}\!c_{\mathfrak a}
      \chi_{\mathfrak a}(p^k)W(N\mathfrak a/D)\right|
       \gg\frac D{\log D}.
\]
Their primary generators have distinct $k$th powers:
equality of the powers would give equality of the prime
ideals by unique factorization, and then equality of their
primary generators. Multiplicativity of the norm gives
$N(p^k)=(Np)^k\le H$. There are consequently
$\gg H^{1/k}/\log D$ distinct such arguments.
Consequently the complete mean square is at least
\[
       \gg\frac{D^2H^{1/k}}{(\log D)^3}.
\]
The ratio of this lower bound to $DH$ grows like
$D^{1-h(1-1/k)}/(\log D)^3$, which tends to infinity.
It even contradicts an estimate $D^{1+\epsilon}H$ for
sufficiently small fixed $\epsilon>0$.
This is an actual counterexample for the enlarged
coefficient class. It is not a counterexample for the
M\"obius coefficients in the supplied papers. It shows
why their special arithmetic identities cannot be replaced
by a statement about arbitrary bounded coefficients.
If a fixed finite-order twist $\nu$ is retained explicitly,
replace $c_{\mathfrak a}$ on prime ideals by
$\overline{\nu(\mathfrak a)}$; the same argument applies.
The counterexample concerns the complete outer set. It does
not assert anything about a restricted set deliberately
excluding $k$th powers. Fixed exclusions remove only
finitely many primes and hence do not change the argument.

The later catalogue entries R107 and R109 retain the same distinction.
Their amplified inverse-moment exponent, with base $D$, inverse
length exponent $1$, no additional prime factors, and outer
range exponent $h<1$, is $(h+5)/6$ after normalizing the inverse
sum by $D^{-1/2}$. Thus their unnormalized bound has exponent
$(11+h)/6$, which exceeds the desired exponent $1+h$ by
$5(1-h)/6$. Their index set is also restricted to
sixth-power-free elements. The more elaborate cofactor bookkeeping
in those entries consequently does not provide the missing
shorter-range mean square.\cite{qrh-catalogue}

There is another useful numerical check. The September marked
inverse moment, with inverse length $D$ and additional prime factors
of total length $D^z$, requires outer range exponent
\[
 h>1+2z,\qquad 3h>2+8z.
\]
If one tries to use those prime factors merely as positive
amplifiers at arguments $p^6$, their product has size
$D^{z/2}$ up to logarithms. The same positivity calculation
then gives the first exponent
\[
 \frac12+\frac{5h}{12}-\frac z2.
\]
The first required inequality already makes this greater than
$11/12+z/3$. Taking $z>0$ therefore does not improve the
unmarked $11/12$ consequence by this particular use of the
marked estimate. The actual $7/8$ proof uses a different
compensated expression and a more detailed estimate; it is
not the direct positive-amplification argument just tested.


% Standalone mathematical sections for inclusion in the reconstruction manuscript.
% Requires amsmath, amssymb, amsthm; theorem, lemma, proposition, corollary,
% and remark environments. All labels and citation keys have the prefix sz-.
